\chapter{Governance and Policy}

Governance is Gert's core differentiator: allowlists, denylists, environment variable blocking,
output redaction, and approval gates. This section defines the v2 governance model, specifying
what policy fields survive from v1, how policy is evaluated at runtime, whether policy is
host-enforced or extension-contributed, and how approval gates integrate with the event model.

Governance features are a first-class design requirement and must not be treated as optional
or defer-to-extension.

% ─────────────────────────────────────────────────────────────────────────────
\section{Governance Primitives}
% ─────────────────────────────────────────────────────────────────────────────

The following primitives form gert's policy vocabulary. Each is defined precisely below, with
its evaluation point in the execution pipeline, its schema representation, and its trace event.

\subsection{Command Allowlist (\texttt{allowed\_commands})}

\textbf{Definition.} A whitelist of executable names (argv[0], without path components) that
may be invoked by \texttt{cli} steps and tool subprocess invocations. If the allowlist is
non-empty, any command whose base name is not in the list is blocked before the subprocess is
forked.

\textbf{Semantics.} Path-stripping is applied before matching: \texttt{/usr/bin/curl} matches
the allowlist entry \texttt{curl}. Glob patterns are \emph{not} supported in v2.0 (exact match
only). The empty list (\texttt{[]}) means all commands are permitted (permissive default).

\textbf{Schema.}
\begin{verbatim}
meta:
  governance:
    allowed_commands: [kubectl, curl, jq, aws]
\end{verbatim}

\textbf{Evaluation point.} Runtime Core, immediately before subprocess fork (after template
evaluation). See §\ref{sec:policy-eval}.

\textbf{Trace event.} \texttt{governance/commandBlocked} — written when a command is denied.
Fields: \texttt{stepId}, \texttt{command}, \texttt{reason: "not\_in\_allowlist"}.

\subsection{Command Denylist (\texttt{denied\_commands})}

\textbf{Definition.} An explicit list of executable names that are unconditionally blocked.
The denylist takes precedence over the allowlist: if a command appears in both, it is denied.

\textbf{Use case.} Operators use the denylist to forbid dangerous commands (\texttt{rm},
\texttt{dd}, \texttt{shutdown}) without needing to enumerate all permitted commands.

\textbf{Schema.}
\begin{verbatim}
meta:
  governance:
    denied_commands: [rm, dd, shutdown, reboot]
\end{verbatim}

\textbf{Evaluation point.} Runtime Core, same checkpoint as allowlist. Denylist check runs
first.

\textbf{Trace event.} \texttt{governance/commandBlocked} — fields: \texttt{stepId},
\texttt{command}, \texttt{reason: "in\_denylist"}.

\subsection{Environment Variable Blocking (\texttt{deny\_env\_vars})}

\textbf{Definition.} A list of glob patterns matched against environment variable names.
Any environment variable whose name matches a pattern is stripped from the subprocess
environment before the command is executed.

\textbf{Rationale.} Prevents accidental or malicious exfiltration of secrets (e.g.,
\texttt{AWS\_SECRET\_ACCESS\_KEY}, \texttt{GITHUB\_TOKEN}) by commands in the runbook.

\textbf{Schema.}
\begin{verbatim}
meta:
  governance:
    deny_env_vars:
      - "*_SECRET*"
      - "*_TOKEN"
      - "AWS_*"
      - "GITHUB_*"
\end{verbatim}

\textbf{Evaluation point.} Runtime Core, immediately before subprocess fork. Applied to the
full process environment (not just explicitly passed vars).

\textbf{Trace event.} \texttt{governance/envVarBlocked} — fields: \texttt{stepId},
\texttt{varNames: [string]} (names of blocked variables, not values).

\subsection{Output Redaction (\texttt{redact})}

\textbf{Definition.} A list of redaction rules. Each rule has a \texttt{pattern} (a RE2
regular expression) and a \texttt{replace} string. After a \texttt{cli} or \texttt{tool}
step produces output, all redaction rules are applied to the captured stdout/stderr before
the output is stored in run variables or written to the trace.

\textbf{Semantics.} Rules are applied in declaration order. The replace string may use RE2
backreferences (\texttt{\$1}). The special replace value \texttt{[REDACTED]} is the
conventional replacement for secrets.

\textbf{Schema.}
\begin{verbatim}
meta:
  governance:
    redact:
      - pattern: "token=[A-Za-z0-9_\\-]{20,}"
        replace: "token=[REDACTED]"
      - pattern: "password: .+"
        replace: "password: [REDACTED]"
      - pattern: "(?i)bearer\\s+[A-Za-z0-9._\\-]+"
        replace: "Bearer [REDACTED]"
\end{verbatim}

\textbf{Evaluation point.} Runtime Core, immediately after subprocess/tool output is captured,
before capture variables are set and before any trace write.

\textbf{Trace event.} \texttt{governance/redactionApplied} — fields: \texttt{stepId},
\texttt{ruleCount} (number of rules that matched, not the patterns themselves, to avoid
logging the pattern that matched sensitive data).

\subsection{Approval Gates}

Approval gates are first-class step-level governance primitives. They are defined on any step
type and pause execution until a minimum number of approvals from specified roles have been
recorded. See §\ref{sec:approval-gate} for the complete approval gate design.

% ─────────────────────────────────────────────────────────────────────────────
\section{Policy Evaluation Model}
\label{sec:policy-eval}
% ─────────────────────────────────────────────────────────────────────────────

Policy primitives are evaluated at specific, well-defined points in the execution pipeline.
The evaluation order within each checkpoint is fixed and must not be reordered by extensions.

\begin{center}
\begin{tabular}{lll}
\textbf{Checkpoint} & \textbf{Primitives Evaluated} & \textbf{Phase} \\
\hline
Parser validation & Schema structural rules & Parse \\
Planner validation & Import cycles, tool existence & Plan \\
Run pre-flight & Extension policy rules & Run init \\
Step pre-flight & Denylist check & Runtime Core \\
Step pre-flight & Allowlist check & Runtime Core \\
Step pre-flight & Approval gate & Runtime Core \\
Subprocess fork & Env var blocking & Runtime Core \\
Output captured & Redaction & Runtime Core \\
Post-step & Contract risk warning & Runtime Core \\
\end{tabular}
\end{center}

\paragraph{Order of precedence.}
Within the pre-flight checkpoint for a \texttt{cli} or \texttt{tool} step, the evaluation
order is:
\begin{enumerate}
  \item Denylist check (hard block, step fails immediately if matched).
  \item Allowlist check (hard block, step fails immediately if not matched and list is non-empty).
  \item Approval gate (soft pause; step is held until approvals arrive).
  \item Env var blocking (applied at fork time, not a failure condition — silently strips).
\end{enumerate}

A policy violation at any hard-block checkpoint writes the corresponding trace event and
returns an error from \texttt{RunHandle.Next()}. The run is placed in the
\texttt{policy\_violation} terminal state. The operator must resolve the policy
(e.g., add the command to the allowlist) and re-run; resumption from a policy violation
is not supported in v2.0.

% ─────────────────────────────────────────────────────────────────────────────
\section{Approval Gate Design}
\label{sec:approval-gate}
% ─────────────────────────────────────────────────────────────────────────────

\subsection{Schema Declaration}

Approval gates are declared in the step's \texttt{approvals} block:

\begin{verbatim}
steps:
  - id: delete_database
    kind: cli
    command: "kubectl delete namespace {{ .targetNamespace }}"
    approvals:
      min: 2
      roles: ["DRI", "change-manager"]
      timeout: 24h
      on_timeout: fail   # or: escalate
      escalate_to: ["senior-sre", "director-of-eng"]
      message: "Deleting {{ .targetNamespace }} is irreversible. Confirm change window."
\end{verbatim}

\subsection{UX Contract}

When the Runtime Core reaches a step with an \texttt{approvals} block, it:
\begin{enumerate}
  \item Writes \texttt{governance/approvalRequired} trace event (stepId, message, roles, min, timeout).
  \item Emits an \texttt{ApprovalRequired} runtime event on the event channel.
  \item Returns from \texttt{Next()} with a \texttt{StepResult} of kind \texttt{PendingApproval}.
        The step does \emph{not} execute yet.
  \item Subsequent calls to \texttt{Next()} on the same step return \texttt{ErrAwaitingApproval}
        until the approval quorum is met.
\end{enumerate}

The adapter (VS Code extension, TUI, or \texttt{gert serve} client) is responsible for
presenting the approval request to authorized users and collecting decisions.

\subsection{Approval Decision Protocol}

An approver submits a decision by calling \texttt{RunHandle.Approve(ctx, ApprovalDecision)}:

\begin{verbatim}
type ApprovalDecision struct {
    StepID   string    // step to approve (must match current paused step)
    Actor    string    // identity of the approver (e.g., "alice@corp.com")
    Role     string    // role the approver is acting in (must be in step.approvals.roles)
    Approved bool      // true = approved, false = rejected
    Comment  string    // optional justification
    At       time.Time // timestamp of decision
}
\end{verbatim}

The runtime validates:
\begin{itemize}
  \item The \texttt{Actor} identity is established (either from \texttt{--as} flag or,
        in a multi-user context, from the adapter's authenticated session).
  \item The \texttt{Role} is in the step's declared \texttt{approvals.roles}.
  \item The same actor may not approve the same gate twice.
  \item If \texttt{Approved = false}, the gate is immediately rejected (any single rejection
        vetos the gate), and the run enters the \texttt{rejected} terminal state.
\end{itemize}

When the approval count reaches \texttt{min}, the runtime writes
\texttt{governance/approvalGranted} to the trace and the next call to \texttt{Next()}
proceeds with normal step execution.

\subsection{Approval Recording in the Trace}

Every approval decision is written to the trace as an immutable record:

\begin{verbatim}
// governance/approvalRequired
{
  "sequence": 14,
  "event": "governance/approvalRequired",
  "stepId": "delete_database",
  "message": "Deleting prod-ns is irreversible. Confirm change window.",
  "roles": ["DRI", "change-manager"],
  "min": 2,
  "timeout": "24h",
  "at": "2026-04-18T10:00:00Z"
}

// governance/approvalDecision (written once per decision)
{
  "sequence": 15,
  "event": "governance/approvalDecision",
  "stepId": "delete_database",
  "actor": "alice@corp.com",
  "role": "DRI",
  "approved": true,
  "comment": "Change window confirmed with customer.",
  "at": "2026-04-18T10:05:00Z"
}

// governance/approvalGranted (written when quorum is met)
{
  "sequence": 16,
  "event": "governance/approvalGranted",
  "stepId": "delete_database",
  "approvals": ["alice@corp.com", "bob@corp.com"],
  "at": "2026-04-18T10:07:00Z"
}
\end{verbatim}

The trace provides a complete, tamper-evident audit record of who approved what and when.

\subsection{Timeout and Escalation}

If \texttt{timeout} is set and the approval quorum is not met within the deadline:
\begin{itemize}
  \item \textbf{\texttt{on\_timeout: fail}:} The run enters the \texttt{timed\_out} terminal state.
        \texttt{governance/approvalTimedOut} is written to the trace.
  \item \textbf{\texttt{on\_timeout: escalate}:} The runtime emits an \texttt{ApprovalEscalated}
        event. If \texttt{escalate\_to} roles are declared, a new \texttt{ApprovalRequired} event
        is emitted targeting those roles. The timeout clock restarts. Escalation is attempted
        once; if the escalated request also times out, the run fails.
\end{itemize}

Timeout enforcement is implemented via a \texttt{context.WithDeadline} passed to the approval
wait loop. The adapter (or the \texttt{gert serve} session timer) is responsible for surfacing
the timeout to the approver.

% ─────────────────────────────────────────────────────────────────────────────
\section{Policy-as-Code Direction}
% ─────────────────────────────────────────────────────────────────────────────

\subsection{v2.0: Structured YAML Policy Blocks}

For v2.0, gert uses \textbf{structured YAML policy blocks} embedded in the runbook's
\texttt{meta.governance} section. This is the approach inherited from v1, extended with the
richer primitives defined in §\ref{sec:approval-gate}.

\textbf{Rationale for v2.0:}
\begin{itemize}
  \item The policy vocabulary is small and closed (six primitives as of v2.0). A full policy
        DSL would be over-engineered for this set.
  \item YAML policy blocks are self-contained: the runbook declares its own governance
        constraints without requiring external policy files or a running policy server.
  \item Structured blocks are straightforward to validate with JSON Schema at authoring time,
        giving operators early feedback in the VS Code extension.
  \item The existing v1 governance model (allowlists, redaction, env blocking) translates
        directly to structured blocks with no loss of expressiveness.
\end{itemize}

\subsection{v2.1+: OPA Integration Hooks}

Dennis's research confirms that Open Policy Agent (OPA) is the CNCF standard for
policy-as-code in operational tooling. For \textbf{v2.1}, gert will introduce OPA
integration at the following evaluation hooks:

\begin{verbatim}
// OPA evaluation hook interface (v2.1 target)
type PolicyEngine interface {
    // EvaluatePreStep is called before each step pre-flight.
    // The policy engine receives the full step context and returns
    // allow/deny with a reason.
    EvaluatePreStep(ctx context.Context, input PolicyInput) (PolicyDecision, error)

    // EvaluatePreRun is called once after plan construction.
    // Used for runbook-level access control (RBAC).
    EvaluatePreRun(ctx context.Context, input PolicyInput) (PolicyDecision, error)
}

type PolicyInput struct {
    Runbook  RunbookMeta
    Step     ResolvedStep
    Actor    string
    Vars     map[string]string
    Contract *contract.Contract
}

type PolicyDecision struct {
    Allow  bool
    Reason string
    // For approval gates: required roles computed by policy engine
    RequiredApprovers []string
}
\end{verbatim}

The \texttt{PolicyEngine} interface will have two implementations in v2.1:
\begin{enumerate}
  \item \textbf{Built-in:} evaluates the structured YAML policy blocks (backward-compatible
        with v2.0 runbooks).
  \item \textbf{OPA:} calls an OPA bundle loaded from a project policy directory
        (\texttt{.gert/policy/*.rego}).
\end{enumerate}

The two implementations compose: the built-in engine runs first; if it allows, the OPA engine
runs second. A denial from either blocks the step.

\textbf{Recommendation:} Do not block the v2.0 launch on OPA integration. Ship the structured
YAML model first, then introduce OPA as an opt-in engine in v2.1. This allows the governance
interface to be designed correctly without the complexity of Rego authoring in the initial release.

% ─────────────────────────────────────────────────────────────────────────────
\section{RBAC Model}
% ─────────────────────────────────────────────────────────────────────────────

\subsection{Identity Establishment}

Actor identity is established via the \texttt{--as <identity>} flag passed at invocation time.
The identity string is an opaque label (conventionally an email address or a role name) recorded
in the trace and used for approval gate matching.

\textbf{v2.0 limitation:} Identity is not authenticated. The \texttt{--as} flag is
self-asserted; gert does not verify that the caller is who they claim to be. Enforcement
of identity authenticity is the responsibility of the surrounding infrastructure (CI/CD
pipeline, VS Code extension SSO, MCP server authentication).

\textbf{v2.1 target:} An \texttt{IdentityProvider} interface will allow adapters to supply
verified identity tokens (e.g., OIDC tokens from VS Code GitHub Copilot authentication),
replacing self-asserted identity for governed environments.

\subsection{Runbook-Level Access Control}

Runbooks may declare a \texttt{requires\_role} field to restrict who may execute them:

\begin{verbatim}
meta:
  governance:
    requires_role: "incident-responder"
    # or: requires_any_role: ["DRI", "incident-responder"]
    # or: requires_all_roles: ["DRI", "change-manager"]
\end{verbatim}

The Runtime Core evaluates this constraint during run initialization (pre-flight, before
any step executes). If the actor's identity does not match, \texttt{governance/accessDenied}
is written to the trace and the run fails immediately.

\subsection{Step-Level Role Restriction}

Individual steps may declare \texttt{requires\_role} to restrict who may trigger that step.
This is distinct from approval gates (which require external approval) — step role restriction
\emph{prevents} a step from executing unless the actor themselves holds the required role.

\begin{verbatim}
steps:
  - id: promote_to_production
    kind: cli
    requires_role: "release-manager"
    command: "kubectl set image deployment/app app={{ .imageTag }}"
\end{verbatim}

\subsection{Separation of Duties}

The approval gate model supports separation of duties: the \texttt{same-actor-cannot-approve}
invariant (defined in §\ref{sec:approval-gate}) ensures that a step author cannot self-approve
their own gate. In environments requiring formal change management, the actor who initiates
the run and the actor(s) who approve high-risk steps must be different.

% ─────────────────────────────────────────────────────────────────────────────
\section{Redaction}
% ─────────────────────────────────────────────────────────────────────────────

\subsection{Pattern-Based Redaction}

The primary redaction mechanism is pattern-based: operators declare RE2 regular expressions
in \texttt{meta.governance.redact}. Redaction is applied eagerly: captured output is
redacted \emph{before} any capture variable assignment or trace write.

\subsection{Explicit Field Marking (v2.1)}

For structured tool outputs (JSON objects returned by MCP tools), v2.1 will introduce explicit
field marking: a tool definition may declare that a returned field contains sensitive data.
The runtime will redact the field value without requiring a regex pattern.

\begin{verbatim}
# In .tool.yaml (v2.1 target)
outputs:
  - name: api_key
    sensitive: true   # value redacted from trace and captures
\end{verbatim}

\subsection{Redaction Guarantees}

\begin{itemize}
  \item Redaction is applied to all output paths: captured variables, trace events, and
        event channel payloads.
  \item Redaction rules are compiled (regex pre-compiled) during run initialization, not
        re-compiled per step.
  \item Redaction is applied unconditionally in all run modes (real, dry-run, replay).
  \item The original (pre-redaction) value is \emph{never} written to disk in any path.
\end{itemize}

% ─────────────────────────────────────────────────────────────────────────────
\section{Audit Trail Requirements}
% ─────────────────────────────────────────────────────────────────────────────

The append-only JSONL trace file is gert's primary audit artifact. Every governance event
is a first-class trace entry. The following governance events are guaranteed to be written:

\begin{center}
\begin{tabular}{ll}
\textbf{Trace Event} & \textbf{Written When} \\
\hline
\texttt{governance/commandBlocked} & A command is blocked by denylist or allowlist \\
\texttt{governance/envVarBlocked} & Env vars are stripped before subprocess fork \\
\texttt{governance/redactionApplied} & Redaction rules matched captured output \\
\texttt{governance/approvalRequired} & A step's approval gate is entered \\
\texttt{governance/approvalDecision} & An approver submits approve or reject \\
\texttt{governance/approvalGranted} & Approval quorum is met \\
\texttt{governance/approvalRejected} & An approver rejects the gate \\
\texttt{governance/approvalTimedOut} & Approval timeout expires \\
\texttt{governance/accessDenied} & RBAC check fails (run or step level) \\
\texttt{governance/contractWarning} & Step contract risk level is critical \\
\texttt{governance/policyViolation} & Generic policy violation (OPA engine, v2.1) \\
\end{tabular}
\end{center}

\subsection{Trace Integrity}

The trace file is protected by:
\begin{itemize}
  \item \textbf{Append-only write discipline:} the trace writer never seeks backward or
        truncates. Each write is followed by \texttt{fsync}.
  \item \textbf{SHA256 evidence hash:} at run completion, the SHA256 of the entire trace
        file is written as the final record and also stored in a sidecar
        \texttt{trace.sha256} file.
  \item \textbf{Monotonic sequence numbers:} every trace event has a \texttt{sequence} field
        (incrementing integer). Gaps in sequence numbers signal corruption or tampering.
\end{itemize}

\subsection{Compliance Alignment}

These audit trail guarantees are designed to satisfy the requirements of:
\begin{itemize}
  \item \textbf{SOC 2 Type II} — Change management controls: every command execution and
        approval is logged with actor identity and timestamp.
  \item \textbf{HIPAA} — Access controls and audit logging: env var blocking and redaction
        prevent PHI from appearing in traces; approval gates enforce separation of duties.
  \item \textbf{ITIL Change Management} — Normal/standard/emergency change paths can be
        modeled via \texttt{requires\_role} and multi-approver gates.
\end{itemize}

% ─────────────────────────────────────────────────────────────────────────────
\section{Extension-Contributed Policy}
% ─────────────────────────────────────────────────────────────────────────────

Extensions may contribute governance rules through the \texttt{ContributedPolicyRules()}
method of the Extension Host (see §02). These rules compose with host policy under the
following invariants:

\begin{enumerate}
  \item \textbf{Additive only.} Extension rules may add restrictions; they may never
        remove or override host-defined policy. An extension cannot remove a command
        from the denylist, relax a required role, or suppress a redaction rule.
  \item \textbf{Host policy takes precedence on conflict.} If an extension rule and a
        host rule address the same primitive and disagree, the more restrictive interpretation
        applies.
  \item \textbf{Extension rules are visible in the trace.} A
        \texttt{governance/extensionPolicyApplied} trace event is written when an
        extension-contributed rule is evaluated, naming the extension and rule ID.
  \item \textbf{Extension rules are not trusted with approval authority.} Extensions may
        not grant approvals on behalf of human approvers. Approval decisions must originate
        from a human actor via the \texttt{RunHandle.Approve()} API.
\end{enumerate}

\subsection{Extension Policy Rule Schema}

\begin{verbatim}
// PolicyRule is contributed by an extension and evaluated at pre-step checkpoint.
type PolicyRule struct {
    ID          string       // unique within the extension's namespace
    Description string
    // Evaluate returns Deny if the rule blocks this step, Allow otherwise.
    // Extensions communicate rules as static declarations (v2.0); dynamic
    // rule evaluation via the extension RPC transport is v2.1.
    StepKinds   []StepKind   // step types this rule applies to (empty = all)
    CommandGlob string       // if set, only applies to commands matching this glob
    Action      PolicyAction // Deny or Warn
    Reason      string       // human-readable reason written to trace event
}
\end{verbatim}

Extension-contributed policy is loaded during the Extension Host's \texttt{Load()} phase,
before the first \texttt{Next()} call. Policy rules are static for the lifetime of a run;
extensions may not add or remove rules after the run has started.
